\documentclass[11pt]{article}
\usepackage{amsmath,amssymb}
\usepackage[margin=1in]{geometry}
\usepackage{hyperref}
\emergencystretch=3em

\title{Headline XIV: the per-stratum posterior expectation in general dimension}
\author{Claude, with Daniel Murfet}
\date{2026-09-07}

\begin{document}
\maketitle
\sloppy

\begin{abstract}
Unit 205 takes the quotient the paper actually cares about. With numerator amplitude $\varphi c$
(observable times density) and denominator amplitude $c>0$, both dressed by the same phase $\xi$,
Headline~XIII gives both leading terms at the same scale, and the denominator coefficient is strictly
positive, so
\[
\frac{\int_{(0,1]^d}\varphi(u)c(u)\,u^he^{-\beta N^2u^{2k}+\beta Nu^k\xi(u)}\,du}
     {\int_{(0,1]^d}c(u)\,u^he^{-\beta N^2u^{2k}+\beta Nu^k\xi(u)}\,du}
\;\longrightarrow\;
\frac{\int\varphi(\pi u)\,c(\pi u)\,J_p(\xi(\pi u))\,w(u)\,du}{\int c(\pi u)\,J_p(\xi(\pi u))\,w(u)\,du}
\]
with $w(u)=\prod_{i\notin J}u_i^{h_i-pk_i}$
(\texttt{headline\_phase\_posterior}): the per-stratum posterior expectation of $\varphi$ concentrates
on the face $u_J=0$ with density $c\,J_p(\xi)\,w$. In the equal-ratio case the face is the corner and
the limit is $\varphi(0)$ (\texttt{headline\_phase\_posterior\_equal}). Zero
\texttt{sorry}/\texttt{axiom}; 9082 jobs.
\end{abstract}

\section{The things formalised}
{\small
\begin{verbatim}
theorem faceProj_mem_closedCube (hu : u ∈ closedCube d) : faceProj h k l u ∈
      closedCube d
theorem continuous_faceProj : Continuous (faceProj h k l)
theorem phaseMoment_pos (hβ) (hp : 0 < p) : 0 < phaseMoment β p a
theorem integral_residualWeight_pos (hk) (hl) (hβ) (hmin) :
    0 < ∫ u in unitBox d, residualWeight h k l u
theorem faceNorm_pos (hk) : 0 < faceNorm h k l
theorem phaseCoeff_pos (hk) (hl) (hβ) (hmin) (hξ) (hc) (hcpos : ∀ x ∈ closedCube d,
      0 < c x) :
    0 < phaseCoeff h k l β ξ c
theorem phase_posterior_tendsto … (hcpos) :
    Tendsto (fun N => (∫ (φ x * c x) * (∏ x^h * quadKernel β (ξ x) (N ∏ x^k)))
        / ∫ c x * (∏ x^h * quadKernel β (ξ x) (N ∏ x^k))) atTop
      (𝓝 (phaseCoeff h k l β ξ (fun x => φ x * c x) / phaseCoeff h k l β ξ c))
theorem headline_phase_posterior … :  … atTop
      (𝓝 ((∫ u in unitBox (m+1), (φ (π u) * c (π u) * phaseMoment β (2l) (ξ (π u)))
      * w u)
          / ∫ u in unitBox (m+1), (c (π u) * phaseMoment β (2l) (ξ (π u))) * w u))
theorem phaseCoeff_equal (hratio) :
    phaseCoeff h k l β ξ η = η 0 * phaseMoment β (2l) (ξ 0) / (m! * ∏ i, k i)
theorem headline_phase_posterior_equal (hratio) (hcpos) : … atTop (𝓝 (φ 0))
\end{verbatim}
}

\section{Mathematics}
Positivity of the denominator coefficient: the face density $G(u)=c(\pi u)J_p(\xi(\pi u))$ is
continuous on the closed cube (\texttt{faceProj} is continuous: each coordinate is either $0$ or a
coordinate projection) and positive there ($c>0$, $J_p>0$ as an integral of a positive integrand
over $(0,\infty)$), so it has a positive minimum $G(x_0)$; then
$\int Gw\ge G(x_0)\int w>0$, where $\int_{(0,1]^d}w>0$ follows from the zero-phase amplitude theorem
at constant amplitude ($\mathrm{faceLeadConst}\cdot\int w=\mathrm{monomialMixedConst}>0$). The
quotient theorem is then \texttt{Tendsto.div} of the two Headline~XIII limits, with the common scale
$N^{-p}(\log N)^{m-1}$ cancelling exactly for $N>1$
(\texttt{div\_div\_div\_cancel\_right₀}). The prefactor $2^{m-1}\cdot2\cdot\mathrm{faceNorm}$ cancels
in the quotient, leaving the ratio of face integrals; in the equal-ratio case the face is the corner
and $\varphi(0)c(0)J_p(\xi(0))/(c(0)J_p(\xi(0)))=\varphi(0)$.

This is the general-$d$, deterministic-phase, single-chart form of the paper's statement that
posterior expectations localise to the exceptional divisor: the observable is read on the face
$u_J=0$, weighted by the density $c$, the dressed normal moment $J_p(\xi)$ of the phase, and the
residual weights of the non-minimising coordinates. It matches the $d=2$ chart posterior limit
$Z[\varphi]/Z[1]\Rightarrow y_{\varphi,00}/y_{1,00}$ of unit~159 (there the amplitudes were
independent Taylor arrays; here $y_{\varphi,00}=\varphi(0)c(0)$, $y_{1,00}=c(0)$).

\section{Lean notes}
\begin{itemize}
\item \texttt{faceProj\_mapsTo} only maps the half-open box into the closed cube; the closed-cube
  version needed here is a two-line lemma (\texttt{faceProj\_mem\_closedCube}).
\item The quotient identity $(a/K)/(b/K)=a/b$ then $(\varphi c J)/(cJ)=\varphi$: rather than
  \texttt{field\_simp} (which left a goal), the three rewrites
  \texttt{div\_div\_div\_cancel\_right₀}, \texttt{mul\_div\_mul\_right}, \texttt{mul\_div\_cancel\_right₀}
  close it exactly.
\item \texttt{IsCompact.exists\_isMinOn} needs a nonemptiness witness; $0\in[0,1]^d$ is
  \texttt{⟨0, fun i \_ => ⟨le\_rfl, zero\_le\_one⟩⟩}.
\item An unused hypothesis (\texttt{hk} in \texttt{phaseCoeff\_equal}) is flagged by the linter;
  drop it from the statement.
\end{itemize}

\section{Status}
Headlines XIII--XIV close the deterministic general-$d$ story at leading order: dressed integral,
cutoff, and posterior quotient. Remaining gaps (report v2 §7): random phases in general dimension,
subleading terms, chart assembly.
\end{document}
